\section{Introduction}
\label{sec:introduction}

Large language models are increasingly used as agents that act in external environments rather than as single-turn text generators. This setting is harder than ordinary preference optimization because credit assignment spans both environment turns and generated tokens, observations are only partial textual renderings of state, and rewards are often sparse.

RAGEN~\citep{wang2025ragen} studies this problem through State-Thinking-Actions-Reward Policy Optimization (StarPO). It combines RAGEN-style prompts, bi-level generalized advantage estimation (GAE), and StarPO-S, a rollout-filtering method that keeps prompts whose sampled trajectories have high reward variance. The original experiments assume server-grade hardware, vLLM-backed rollout, fp32 AdamW optimizer states, and multi-seed runs.

This project asks a narrower question: how much of the RAGEN training pipeline can be reproduced on a single consumer GPU, and what changes when the update algorithm is switched from PPO with a critic to critic-free GRPO? The implementation preserves the main RAGEN contracts, rewrites the system into modular components, and makes the hardware concessions explicit.

The paper first reviews the necessary background on LLM-agent RL, PPO/GRPO, and StarPO-S filtering, then describes the implementation method, experimental benchmark, results, and limitations. Additional curves, detailed tables, implementation settings, and training commands are provided in the appendix. The contribution of this report is summarized as follows:

\begin{itemize}[leftmargin=*]
    \item \textbf{Consumer-hardware reproduction.} \method{} rewrites the pipeline into five modules: environments, agents, RL algorithms, rollout/training, and evaluation. It adds an 8 GB VRAM optimization stack, including AdamW8bit with fp32 embedding states, gradient checkpointing, bf16 weights, micro-batching, KV-cache restoration for generation, batched rollout, and alive-only batch shrinking. These changes make 200-step 0.5B-model PPO and GRPO runs feasible on a single RTX 4070.
    \item \textbf{Algorithmic boundary of variance filtering.} The project completes a \(2 \times 3\) experiment matrix over algorithm \(\in\) \{PPO, GRPO\} and variance-filter ratio \(\in \{1.0, 0.5, 0.25\}\). PPO directionally reproduces the RAGEN story: vanilla training collapses, moderate filtering mitigates collapse, and aggressive filtering becomes stagnant. GRPO reverses the pattern: vanilla GRPO is the strongest configuration. This suggests that variance filtering is not algorithm-agnostic; it helps PPO partly through critic-data selection, but does not transfer directly to actor-only GRPO.
\end{itemize}
